% BÀI TOÁN HÌNH HỌC - BAI_HINH_29 (TRÍCH ĐỀ 29)
\begin{ex}
Cho nửa đường tròn tâm O đường kính AB = 2R. Từ A vẽ tia tiếp tuyến Ax cùng phía với nửa đường tròn đối với AB. Từ điểm M trên Ax kẻ tiếp tuyến thứ hai MC với nửa đường tròn (C là tiếp điểm). MB cắt nửa đường tròn (O) tại D (D khác B).
	\begin{enumerate}
		\item[a)] Chứng minh A, M, C, O cùng nằm trên một đường tròn.
		\item[b)] Vẽ CH vuông góc với AB ($H \in AB$). Chứng minh rằng: $MA^2 = MD \cdot MB$ và MB đi qua trung điểm của CH.
	\end{enumerate}
\end{ex}
